\documentclass[11pt,a4paper]{article}

% ── Codificación, fuentes y lenguaje ───────────────────────────────────────────
\usepackage[utf8]{inputenc}
\usepackage[T1]{fontenc}
\usepackage[default,scale=0.95]{sourcesanspro}
\renewcommand{\familydefault}{\sfdefault}
\usepackage[spanish,es-noshorthands,es-tabla]{babel}

% ── Geometría y espaciado ─────────────────────────────────────────────────────
\usepackage[top=2.2cm, bottom=2.5cm, left=2.2cm, right=2.2cm, headheight=15pt]{geometry}
\usepackage{setspace}
\setstretch{1.18}
\usepackage{parskip}

% ── Matemáticas y unidades ────────────────────────────────────────────────────
\usepackage{amsmath}
\usepackage{amssymb}
\usepackage{siunitx}

% ── Circuitos ─────────────────────────────────────────────────────────────────
\usepackage[european, straightvoltages]{circuitikz}

% ── Tablas ────────────────────────────────────────────────────────────────────
\usepackage{booktabs}
\usepackage{tabularx}
\usepackage{array}
\usepackage{longtable}
\usepackage{colortbl}
\usepackage{enumitem}

% ── Colores, cajas e iconos ───────────────────────────────────────────────────
\usepackage[dvipsnames]{xcolor}
\usepackage{tcolorbox}
\tcbuselibrary{skins, breakable}
\usepackage{fontawesome5}
\usepackage{graphicx}

% ── Listados de código (dark-mode infográfico) ────────────────────────────────
%   Estilo base: fondo oscuro con números de línea. Los bloques lstlisting se
%   envuelven automáticamente en una caja tcolorbox redondeada.
\usepackage{listings}

% ── Secciones con barra de color (infografía) ────────────────────────────────
\usepackage{titlesec}
\titleformat{\section}
  {\normalfont\LARGE\bfseries\color{azulNoche}}
  {\colorbox{cyanNeon}{\color{fondoOscuro}\thesection}}{0.7em}{}
  [\vspace{2pt}\color{cyanNeon}\rule{\linewidth}{1.8pt}]
\titlespacing*{\section}{0pt}{24pt}{10pt}
\titleformat{\subsection}
  {\normalfont\large\bfseries\color{azulNoche}}
  {\textcolor{cyanNeon}{\thesubsection}}{0.7em}{}
  [\vspace{1pt}\color{grisLinea}\rule{\linewidth}{0.6pt}]
\titlespacing*{\subsection}{0pt}{14pt}{6pt}
\titleformat{\subsubsection}
  {\normalfont\normalsize\bfseries\color{azulNoche}}
  {\textcolor{cyanNeon}{\thesubsubsection}}{0.7em}{}
\titlespacing*{\subsubsection}{0pt}{10pt}{4pt}

% ── Cabeceras y pies de página ────────────────────────────────────────────────
\usepackage{fancyhdr}
\usepackage{lastpage}
\pagestyle{fancy}
\fancyhf{}
\renewcommand{\headrulewidth}{0pt}
\renewcommand{\footrulewidth}{0pt}
\fancyfoot[C]{%
  \begin{minipage}{\linewidth}
    \vspace{2pt}
    \color{cyanNeon}\rule{\linewidth}{1.2pt}\\[2pt]
    \centering\color{grisTexto}\footnotesize
    Página \thepage\ de \pageref{LastPage}
  \end{minipage}%
}

% ── Hiperenlaces ──────────────────────────────────────────────────────────────
\usepackage[colorlinks=true, linkcolor=azulNoche, urlcolor=cyanNeon]{hyperref}
\sloppy
\emergencystretch 3em

% ── Paleta de colores — tecnológico dark-mode ──────────────────────────────────
\definecolor{fondoOscuro}{HTML}{0D1B2A}%             Dark navy — fondo banda título
\definecolor{verdeTurquesa}{HTML}{004D40}%           Verde turquesa — bandas de marca
\definecolor{azulNoche}{HTML}{1B3A6B}%               Midnight blue — secciones, marcos
\definecolor{azulMedio}{HTML}{2A5298}%               Azul medio — variante de acento
\definecolor{cyanNeon}{HTML}{00C8E0}%                Cyan eléctrico — acentos primarios
\definecolor{verdeSignal}{HTML}{00B887}%             Verde señal — fortalezas, éxito
\definecolor{naranjaVivo}{HTML}{FF7043}%             Naranja vivo — mejoras, advertencias
\definecolor{rojoAlerta}{HTML}{E53935}%              Rojo alerta — errores
\definecolor{grisPapel}{HTML}{F4F6FA}%               Fondo tarjetas claras
\definecolor{grisLinea}{HTML}{C8D0E0}%               Bordes sutiles
\definecolor{grisTexto}{HTML}{6B7A99}%               Texto secundario
\definecolor{amarilloNota}{HTML}{FFB300}%            Notas / advertencias doradas
\definecolor{codigoFondo}{HTML}{1E2D3D}%             Fondo del bloque de código
\definecolor{codigoComentario}{HTML}{7A8FA6}%        Comentarios
\definecolor{codigoCadena}{HTML}{FF9D5C}%            Cadenas / strings
\definecolor{codigoKeyword}{HTML}{00C8E0}%           Palabras clave
\definecolor{codigoNumero}{HTML}{00E5A0}%            Números

% ── Banda de título: portada infográfica ──────────────────────────────────────
%   Diseño en 2 capas: banda de color (por defecto fondo oscuro) + franja de
%   acento cyan abajo. Uso: \bandaTitulo[color]{título}{subtítulo}.
%   El icono se puede cambiar con \renewcommand{\iconoBanda}{...} (FontAwesome).
\newcommand{\iconoBanda}{microchip}
\newcommand{\bandaTitulo}[3][fondoOscuro]{%
  \begin{tcolorbox}[
    enhanced,
    colback=#1, colframe=#1,
    boxrule=0pt, arc=8pt, outer arc=8pt,
    width=\linewidth,
    top=18pt, bottom=6pt, left=14pt, right=14pt,
    borderline south={3.5pt}{0pt}{cyanNeon},
  ]
    \begin{minipage}[c]{0.07\linewidth}
      \centering
      \textcolor{cyanNeon}{\fontsize{30}{30}\selectfont\faIcon{\iconoBanda}}
    \end{minipage}%
    \hspace{8pt}%
    \begin{minipage}[c]{0.88\linewidth}
      {\fontsize{20}{24}\selectfont\bfseries\color{white}#2}\par\vspace{5pt}%
      \textcolor{cyanNeon!80!white}{\small\faIcon{angle-right}~#3}%
    \end{minipage}
    \vspace{4pt}
  \end{tcolorbox}%
  \vspace{8pt}%
}

% ── Tarjetas de datos (infografía) ────────────────────────────────────────────
\tcbset{
  tarjetaDato/.style={
    enhanced, breakable,
    arc=6pt, outer arc=6pt,
    colback=grisPapel, colframe=azulNoche,
    boxrule=1pt,
    fonttitle=\bfseries\small, coltitle=white,
    attach boxed title to top left={yshift=-3mm, xshift=6mm},
    boxed title style={
      colback=azulNoche, colframe=azulNoche,
      arc=4pt, boxrule=0pt,
    },
    drop fuzzy shadow=azulNoche!30!white,
    top=8pt, bottom=8pt, left=10pt, right=10pt,
  },
  tarjetaIcono/.style={
    enhanced, breakable,
    arc=6pt, outer arc=6pt,
    colback=white, colframe=grisLinea, boxrule=0.8pt,
    fonttitle=\bfseries\small, coltitle=azulNoche,
    borderline west={3pt}{0pt}{cyanNeon},
    drop fuzzy shadow=azulNoche!25!white,
    top=6pt, bottom=6pt, left=10pt, right=10pt,
  },
  cajaTitulo/.style={
    enhanced, breakable,
    arc=6pt, outer arc=6pt,
    colback=azulNoche!10!grisPapel, colframe=azulNoche,
    boxrule=1pt,
    fonttitle=\bfseries, coltitle=white,
    attach boxed title to top left={yshift=-3mm, xshift=6mm},
    boxed title style={
      colback=azulNoche, colframe=azulNoche,
      arc=4pt, boxrule=0pt,
    },
    drop fuzzy shadow=azulNoche!25!white,
    top=8pt, bottom=8pt, left=10pt, right=10pt,
  },
  cajaContenido/.style={
    enhanced, breakable,
    arc=5pt, outer arc=5pt,
    colback=grisPapel, colframe=grisLinea,
    boxrule=0.8pt,
    fonttitle=\bfseries\small, coltitle=azulNoche,
    top=6pt, bottom=6pt, left=10pt, right=10pt,
  },
  cajaFortaleza/.style={
    enhanced, breakable,
    arc=6pt, outer arc=6pt,
    colback=verdeSignal!8!white, colframe=verdeSignal,
    boxrule=1pt,
    borderline west={4pt}{0pt}{verdeSignal},
    fonttitle=\bfseries\small, coltitle=verdeSignal!70!black,
    drop fuzzy shadow=verdeSignal!20!white,
    top=6pt, bottom=6pt, left=10pt, right=10pt,
  },
  cajaMejora/.style={
    enhanced, breakable,
    arc=6pt, outer arc=6pt,
    colback=naranjaVivo!8!white, colframe=naranjaVivo,
    boxrule=1pt,
    borderline west={4pt}{0pt}{naranjaVivo},
    fonttitle=\bfseries\small, coltitle=naranjaVivo!80!black,
    drop fuzzy shadow=naranjaVivo!20!white,
    top=6pt, bottom=6pt, left=10pt, right=10pt,
  },
  cajaRecomendacion/.style={
    enhanced, breakable,
    arc=6pt, outer arc=6pt,
    colback=cyanNeon!6!white, colframe=cyanNeon!60!azulNoche,
    boxrule=1pt,
    borderline west={4pt}{0pt}{cyanNeon},
    fonttitle=\bfseries\small, coltitle=azulNoche,
    drop fuzzy shadow=azulNoche!20!white,
    top=6pt, bottom=6pt, left=10pt, right=10pt,
  },
}

% ── Ícono + texto (encabezados de sección y elementos) ────────────────────────
\newcommand{\iconotexto}[2]{\textcolor{cyanNeon}{\faIcon{#1}}~#2}

% ── Listados de código: estilo dark + auto-envoltura ─────────────────────────
\lstdefinestyle{estiloCodigo}{
    backgroundcolor=\color{codigoFondo},
    basicstyle=\ttfamily\small\color{white},
    commentstyle=\color{codigoComentario}\itshape,
    stringstyle=\color{codigoCadena},
    keywordstyle=\color{codigoKeyword}\bfseries,
    numberstyle=\tiny\color{codigoComentario},
    numbers=left,
    numbersep=10pt,
    showstringspaces=false,
    breaklines=true,
    breakatwhitespace=false,
    tabsize=4,
    frame=none,
    captionpos=b,
    aboveskip=6pt,
    belowskip=4pt,
    xleftmargin=14pt,
    framexleftmargin=14pt,
    literate=
      {á}{{\'a}}1 {é}{{\'e}}1 {í}{{\'i}}1 {ó}{{\'o}}1 {ú}{{\'u}}1
      {Á}{{\'A}}1 {É}{{\'E}}1 {Í}{{\'I}}1 {Ó}{{\'O}}1 {Ú}{{\'U}}1
      {ñ}{{\~n}}1 {Ñ}{{\~N}}1 {ü}{{\"u}}1 {Ü}{{\"U}}1
      {¿}{{?`}}1 {¡}{{!`}}1
}
\lstdefinestyle{estiloPython}{
    style=estiloCodigo,
    language=Python,
}
\lstset{style=estiloCodigo}
\BeforeBeginEnvironment{lstlisting}{%
  \begin{tcolorbox}[
    enhanced, breakable,
    arc=6pt, outer arc=6pt,
    colback=codigoFondo, colframe=azulNoche,
    boxrule=1pt,
    drop fuzzy shadow=azulNoche!25!white,
    top=2pt, bottom=2pt, left=0pt, right=0pt,
  ]%
}
\AfterEndEnvironment{lstlisting}{\end{tcolorbox}}

% ── Cajas de contenido temáticas ──────────────────────────────────────────────
\newtcolorbox{cajaRecuerda}{
    enhanced, breakable,
    arc=6pt, outer arc=6pt,
    colback=azulNoche!10!grisPapel,
    colframe=azulNoche, boxrule=1pt,
    borderline west={4pt}{0pt}{cyanNeon},
    fonttitle=\bfseries\small\color{white},
    title={\faIcon{info-circle}~Recuerda},
    attach boxed title to top left={yshift=-3mm, xshift=6mm},
    boxed title style={colback=azulNoche, arc=4pt, boxrule=0pt},
    drop fuzzy shadow=azulNoche!25!white,
    left=10pt, right=10pt, top=8pt, bottom=8pt,
}

\newtcolorbox{cajaConcepto}{
    enhanced, breakable,
    arc=6pt, outer arc=6pt,
    colback=verdeSignal!8!white,
    colframe=verdeSignal, boxrule=1pt,
    borderline west={4pt}{0pt}{verdeSignal},
    fonttitle=\bfseries\small\color{white},
    title={\faIcon{lightbulb}~Concepto clave},
    attach boxed title to top left={yshift=-3mm, xshift=6mm},
    boxed title style={colback=verdeSignal!80!black, arc=4pt, boxrule=0pt},
    drop fuzzy shadow=verdeSignal!20!white,
    left=10pt, right=10pt, top=8pt, bottom=8pt,
}

\newtcolorbox{cajaEjemplo}{
    enhanced, breakable,
    arc=6pt, outer arc=6pt,
    colback=cyanNeon!5!white,
    colframe=cyanNeon!70!azulNoche, boxrule=1pt,
    borderline west={4pt}{0pt}{cyanNeon},
    fonttitle=\bfseries\small\color{fondoOscuro},
    title={\faIcon{code}~Ejemplo},
    attach boxed title to top left={yshift=-3mm, xshift=6mm},
    boxed title style={colback=cyanNeon, arc=4pt, boxrule=0pt},
    drop fuzzy shadow=azulNoche!20!white,
    left=10pt, right=10pt, top=8pt, bottom=8pt,
}

\newtcolorbox{cajaEjercicio}{
    enhanced, breakable,
    arc=6pt, outer arc=6pt,
    colback=azulMedio!7!white,
    colframe=azulMedio, boxrule=1pt,
    borderline west={4pt}{0pt}{azulMedio},
    fonttitle=\bfseries\small\color{white},
    title={\faIcon{pencil-alt}~Ejercicio},
    attach boxed title to top left={yshift=-3mm, xshift=6mm},
    boxed title style={colback=azulMedio, arc=4pt, boxrule=0pt},
    drop fuzzy shadow=azulNoche!20!white,
    left=10pt, right=10pt, top=8pt, bottom=8pt,
}

\newtcolorbox{cajaReto}{
    enhanced, breakable,
    arc=6pt, outer arc=6pt,
    colback=naranjaVivo!6!white,
    colframe=naranjaVivo, boxrule=1pt,
    borderline west={4pt}{0pt}{naranjaVivo},
    fonttitle=\bfseries\small\color{white},
    title={\faIcon{fire}~Reto},
    attach boxed title to top left={yshift=-3mm, xshift=6mm},
    boxed title style={colback=naranjaVivo, arc=4pt, boxrule=0pt},
    drop fuzzy shadow=naranjaVivo!20!white,
    left=10pt, right=10pt, top=8pt, bottom=8pt,
}

% ── Indicadores de dificultad ─────────────────────────────────────────────────
\newcommand{\facil}{\textcolor{verdeSignal}{\faIcon{circle}\,\textbf{Fácil}}}
\newcommand{\intermedio}{\textcolor{naranjaVivo}{\faIcon{circle}\,\textbf{Intermedio}}}
\newcommand{\dificil}{\textcolor{rojoAlerta}{\faIcon{circle}\,\textbf{Difícil}}}

% ─────────────────────────────────────────────────────────────────────────────
\renewcommand{\iconoBanda}{flask}
\fancyhead[L]{\footnotesize\color{grisTexto}\faIcon{flask}~LABORATORIO\_I\_FISICA}
\fancyhead[R]{\footnotesize\color{grisTexto}Plan semanal}
\hypersetup{pdftitle={Laboratorio I de Física},pdfauthor={Prof. César Rodríguez}
  pdfsubject={LABORATORIO\_I\_FISICA}}

\begin{document}
\bandaTitulo[fondoOscuro]{Plan semanal}{Laboratorio I de Física --- 10 semanas}

\begin{center}\small\color{grisTexto}Prof. César Rodríguez \textbar{} Universidad de Oriente, Núcleo de Sucre, Departamento de Física \textbar{} 6 de noviembre de 2026\end{center}

Cada semana se entrega en GitHub, dentro de la carpeta semana\_NN, antes del corte del domingo a las 23:59. El docente revisa esa entrega al finalizar la semana y devuelve observaciones antes de la siguiente.

\begin{tcolorbox}[cajaContenido,title={Semana 1 --- La medición, la apreciación y el error}]
\textbf{Fase 1} --- Medición, error y análisis gráfico

\textbf{Objetivo.} Distinguir apreciación, exactitud y precisión, expresar una lectura con su incertidumbre y decidir cuántas cifras significativas tiene sentido reportar.

\textbf{Contenidos}

\begin{itemize}[leftmargin=1.4em,itemsep=2pt,topsep=3pt]
  \item Magnitud, unidad y patrón; magnitudes directas e indirectas.
  \item Apreciación de un instrumento de escala dividida y de un instrumento digital.
  \item Exactitud frente a precisión; errores sistemáticos y aleatorios.
  \item Error absoluto, relativo y porcentual.
  \item Valor medio, desviación estándar y desviación estándar de la media.
  \item Incertidumbre tipo A y tipo B; incertidumbre combinada.
  \item Cifras significativas y reglas de redondeo.
\end{itemize}

\textbf{Herramientas:} Regla milimetrada, Calibrador o pie de rey, Tornillo micrométrico, Cronómetro digital, Calculadora.

\textbf{Taller.} \emph{Medición, apreciación y expresión de resultados}. Cada grupo recibe tres instrumentos con distinta apreciación y un juego de cuerpos. Debe medir la misma longitud con los tres, repetir diez veces la lectura con el instrumento de menor apreciación y decidir cuál de las tres expresiones del resultado es defendible.

\textbf{Entregable:} Taller resuelto con la tabla de diez lecturas, el cálculo completo de la media y las incertidumbres, la expresión final del resultado y un párrafo que justifique el instrumento elegido.

\textbf{Instrumento de evaluación:} Taller formativo con entrega.
\end{tcolorbox}

\begin{tcolorbox}[cajaContenido,title={Semana 2 --- Propagación de incertidumbres y notación científica}]
\textbf{Fase 1} --- Medición, error y análisis gráfico

\textbf{Objetivo.} Propagar la incertidumbre desde las magnitudes medidas hasta una magnitud calculada, y expresar el resultado en notación científica con las cifras significativas que la incertidumbre permite.

\textbf{Contenidos}

\begin{itemize}[leftmargin=1.4em,itemsep=2pt,topsep=3pt]
  \item Regla de la suma y de la diferencia para incertidumbres absolutas.
  \item Regla del producto y del cociente para incertidumbres relativas.
  \item Propagación general por derivadas parciales.
  \item Potencias, raíces y factores constantes.
  \item Notación científica y prefijos del SI.
  \item Cifras significativas en operaciones y redondeo final.
\end{itemize}

\textbf{Herramientas:} Calculadora científica, Calibrador, Probeta graduada, Hoja de cálculo o Python.

\textbf{Taller.} \emph{Del objeto al resultado: propagar sin inventar cifras}. Cada grupo mide el diámetro y la altura de un cilindro con el calibrador y calcula su volumen por la fórmula geométrica. Debe reportar el volumen con su incertidumbre, no con todas las cifras que arroja la calculadora, y comprobar el resultado por desplazamiento de agua.

\textbf{Entregable:} Taller resuelto con las mediciones, la propagación paso a paso, ambos volúmenes expresados en notación científica y la conclusión sobre si los métodos coinciden dentro de la incertidumbre.

\textbf{Instrumento de evaluación:} Taller formativo con entrega.
\end{tcolorbox}

\begin{tcolorbox}[cajaContenido,title={Semana 3 --- Graficación, linealización y lectura de una recta}]
\textbf{Fase 1} --- Medición, error y análisis gráfico

\textbf{Objetivo.} Construir una gráfica con ejes rotulados y unidades, linealizar un modelo físico y extraer de la recta la pendiente y el intercepto con su significado físico.

\textbf{Contenidos}

\begin{itemize}[leftmargin=1.4em,itemsep=2pt,topsep=3pt]
  \item Variables dependiente e independiente; elección de escalas y origen.
  \item Rotulación de ejes, unidades y título de figura.
  \item Barras de incertidumbre y su relación con el tamaño del punto.
  \item Modelos linealizables: potencial, exponencial y proporcional.
  \item Pendiente e intercepto con unidades y significado físico.
  \item Criterio de bondad visual: dispersión de los puntos y tendencia.
\end{itemize}

\textbf{Herramientas:} Papel milimetrado, Regla, Datos de la práctica del péndulo, Python con Matplotlib.

\textbf{Taller.} \emph{De la curva a la recta: linealizar para entender}. Cada grupo recibe una tabla de periodo y longitud de un péndulo. Debe graficar los datos crudos, descubrir que la relación no es lineal, linealizarla elevando al cuadrado el periodo y leer de la recta la aceleración de la gravedad.

\textbf{Entregable:} Taller resuelto con las dos gráficas, la recta trazada, el cálculo de la pendiente con unidades, el valor de $g$ y la interpretación del intercepto.

\textbf{Instrumento de evaluación:} Taller formativo con entrega.
\end{tcolorbox}

\begin{tcolorbox}[cajaContenido,title={Semana 4 --- Mínimos cuadrados y estructura del informe}]
\textbf{Fase 1} --- Medición, error y análisis gráfico

\textbf{Objetivo.} Ajustar una recta por mínimos cuadrados, reportar la pendiente, el intercepto y su bondad, y estructurar un informe de laboratorio que otro pueda reproducir.

\textbf{Contenidos}

\begin{itemize}[leftmargin=1.4em,itemsep=2pt,topsep=3pt]
  \item Criterio de mínimos cuadrados y diferencia con el trazo a ojo.
  \item Fórmulas de pendiente e intercepto para datos con x sin incertidumbre.
  \item Coeficiente de determinación y análisis de residuales.
  \item Incertidumbre de la pendiente y del intercepto.
  \item Estructura del informe: objetivo, fundamento, montaje, datos, análisis, conclusiones.
  \item Trazabilidad: cada cifra del texto sale de una tabla o de un script.
\end{itemize}

\textbf{Herramientas:} Hoja de cálculo o Python con NumPy, Datos del resorte de la semana 8, Calculadora.

\textbf{Taller.} \emph{Ajustar, juzgar y escribir: el cierre de la fase uno}. Cada grupo recibe una tabla de masa colgada y estiramiento de un resorte. Debe ajustar la recta por mínimos cuadrados, decidir si el resorte es lineal en el rango medido y redactar un informe corto con la estructura completa del curso.

\textbf{Entregable:} Taller integrador con la tabla completa, el ajuste por mínimos cuadrados, la gráfica con residuales, la constante elástica y un informe corto de dos páginas con la estructura del curso.

\textbf{Instrumento de evaluación:} Taller integrador con entrega.
\end{tcolorbox}

\begin{tcolorbox}[cajaContenido,title={Semana 5 --- Volumen de un sólido por dos métodos}]
\textbf{Fase 2} --- Propiedades de la materia y equilibrio

\textbf{Objetivo.} Determinar el volumen de un sólido regular por medición geométrica y por desplazamiento de agua, propagar la incertidumbre de ambos métodos y decidir si sus resultados coinciden dentro del error experimental.

\textbf{Contenidos}

\begin{itemize}[leftmargin=1.4em,itemsep=2pt,topsep=3pt]
  \item Volumen de cuerpos regulares: prisma, cilindro y esfera.
  \item Medición de dimensiones con calibrador y tornillo micrométrico.
  \item Desplazamiento de líquido y lectura del menisco.
  \item Propagación de incertidumbres en productos y potencias.
  \item Comparación de dos métodos independientes.
\end{itemize}

\textbf{Herramientas:} Calibrador y tornillo micrométrico, Probeta graduada de 100 mL, Agua destilada, jeringa y paño, Cuerpos regulares de distintos materiales, Hilo fino.

\textbf{Práctica.} \emph{Volumen de un sólido por dos métodos}. Cada grupo recibe tres sólidos regulares. Debe medir sus dimensiones con el instrumento adecuado, calcular el volumen geométrico, medirlo de nuevo por desplazamiento de agua y explicar la discrepancia entre ambos resultados con su presupuesto de incertidumbres.

\textbf{Entregable:} Informe con la tabla de datos crudos, ambos volúmenes con su incertidumbre, la comparación entre métodos y la explicación de la discrepancia.

\textbf{Instrumento de evaluación:} Práctica con informe.
\end{tcolorbox}

\begin{tcolorbox}[cajaContenido,title={Semana 6 --- Densidad de un cuerpo sólido}]
\textbf{Fase 2} --- Propiedades de la materia y equilibrio

\textbf{Objetivo.} Determinar la densidad de un sólido por la relación masa-volumen y por el principio de Arquímedes, comparar ambos resultados y reconocer el material a partir de la tabla de densidades.

\textbf{Contenidos}

\begin{itemize}[leftmargin=1.4em,itemsep=2pt,topsep=3pt]
  \item Densidad como propiedad intensiva de la materia.
  \item Medición de masa con balanza y de volumen por desplazamiento.
  \item Principio de Arquímedes y empuje.
  \item Peso aparente y determinación de densidad sin medir el volumen.
  \item Propagación de incertidumbres en un cociente.
  \item Dependencia de la densidad del agua con la temperatura.
\end{itemize}

\textbf{Herramientas:} Balanza de precisión de 0,01 g, Calibrador y probeta graduada, Vaso de precipitados y agua destilada, Termómetro digital, Hilo fino y paño, Sólidos metálicos de distintos materiales.

\textbf{Práctica.} \emph{Densidad de un cuerpo sólido}. Cada grupo recibe varios sólidos metálicos sin identificar. Debe obtener la densidad de cada uno por la relación masa-volumen y por pesada hidrostática, comparar los dos valores y proponer de qué material se trata usando la tabla de densidades.

\textbf{Entregable:} Informe con la tabla de datos, las densidades por ambos métodos con su incertidumbre, la identificación propuesta del material y el análisis de las fuentes de error.

\textbf{Instrumento de evaluación:} Práctica con informe.
\end{tcolorbox}

\begin{tcolorbox}[cajaContenido,title={Semana 7 --- Centro de masa de un sistema de dos cuerpos}]
\textbf{Fase 2} --- Propiedades de la materia y equilibrio

\textbf{Objetivo.} Localizar experimentalmente el centro de masa de un sistema formado por una varilla y dos masas, compararlo con el valor calculado a partir de las masas y las posiciones, y analizar el efecto de la masa de la varilla.

\textbf{Contenidos}

\begin{itemize}[leftmargin=1.4em,itemsep=2pt,topsep=3pt]
  \item Centro de masa de un sistema discreto.
  \item Centro de masa de una varilla uniforme.
  \item Equilibrio de un cuerpo extenso y momento de una fuerza.
  \item Medición de posiciones respecto de un origen arbitrario.
  \item Propagación de incertidumbres en un promedio ponderado.
\end{itemize}

\textbf{Herramientas:} Varilla metálica de longitud conocida, Dos masas con tornillo de fijación, Soporte de cuchilla o punto de apoyo, Balanza digital, Regla milimetrada y calibrador, Nivel de burbuja.

\textbf{Práctica.} \emph{Centro de masa de un sistema de dos cuerpos}. Cada grupo monta dos masas sobre una varilla, localiza el punto de equilibrio y lo compara con el centro de masa calculado. Debe repetir el montaje con tres separaciones distintas y decidir si la masa de la varilla puede despreciarse.

\textbf{Entregable:} Informe con la tabla de los tres montajes, los centros de masa calculados con y sin varilla, la comparación con el medido y la conclusión sobre la validez de despreciar la masa de la varilla.

\textbf{Instrumento de evaluación:} Práctica con informe.
\end{tcolorbox}

\begin{tcolorbox}[cajaContenido,title={Semana 8 --- Ley de Hooke y constante elástica}]
\textbf{Fase 3} --- Elasticidad, oscilaciones y calor

\textbf{Objetivo.} Verificar la proporcionalidad entre la fuerza aplicada y el estiramiento de un resorte, y determinar su constante elástica a partir de la pendiente de un ajuste por mínimos cuadrados.

\textbf{Contenidos}

\begin{itemize}[leftmargin=1.4em,itemsep=2pt,topsep=3pt]
  \item Ley de Hooke y régimen elástico.
  \item Fuerza aplicada como peso de masas calibradas.
  \item Estiramiento respecto de la posición de equilibrio sin carga.
  \item Ajuste lineal de fuerza frente a estiramiento.
  \item Constante elástica, unidades y significado físico.
  \item Límite elástico y desviación de la linealidad.
\end{itemize}

\textbf{Herramientas:} Soporte universal con pinza, Resorte helicoidal, Juego de masas calibradas, Regla milimetrada o cinta métrica, Balanza digital, Calculadora o Python.

\textbf{Práctica.} \emph{Ley de Hooke y constante elástica}. Cada grupo cuelga masas crecientes de un resorte, mide el estiramiento y ajusta la recta fuerza-estiramiento. Debe reportar la constante elástica con su incertidumbre y señalar hasta qué masa el resorte sigue siendo lineal.

\textbf{Entregable:} Informe con la tabla de datos, la gráfica con el ajuste y los residuales, la constante elástica con su incertidumbre y el rango de validez de la ley de Hooke.

\textbf{Instrumento de evaluación:} Práctica con informe.
\end{tcolorbox}

\begin{tcolorbox}[cajaContenido,title={Semana 9 --- Péndulo simple y determinación de la gravedad}]
\textbf{Fase 3} --- Elasticidad, oscilaciones y calor

\textbf{Objetivo.} Determinar la aceleración de la gravedad a partir de la linealización del periodo del péndulo simple y estimar la corrección de longitud asociada al centro de masa del cuerpo oscilante.

\textbf{Contenidos}

\begin{itemize}[leftmargin=1.4em,itemsep=2pt,topsep=3pt]
  \item Péndulo simple y aproximación de ángulo pequeño.
  \item Periodo, longitud efectiva y centro de masa.
  \item Linealización de periodo al cuadrado frente a longitud.
  \item Ajuste por mínimos cuadrados y obtención de la gravedad.
  \item Interpretación del intercepto como corrección de longitud.
  \item Errores dominantes: ángulo, cronometraje y medición de longitud.
\end{itemize}

\textbf{Herramientas:} Soporte con hilo inextensible, Esfera metálica de masa conocida, Cinta métrica y calibrador, Cronómetro digital, Transportador, Calculadora o Python.

\textbf{Práctica.} \emph{Péndulo simple y determinación de la gravedad}. Cada grupo recibe una tabla de veintiuna longitudes de hilo con el tiempo de cinco oscilaciones. Debe linealizar el modelo, ajustar la recta, obtener la gravedad y explicar por qué el intercepto no es cero. Los datos provienen de una medición real de laboratorio.

\textbf{Entregable:} Informe con la tabla completa, la gráfica linealizada, el ajuste, el valor de $g$ con su incertidumbre, la corrección de longitud y el análisis de los datos anómalos.

\textbf{Instrumento de evaluación:} Práctica con informe y datos reales.
\end{tcolorbox}

\begin{tcolorbox}[cajaContenido,title={Semana 10 --- Enfriamiento de un cuerpo sólido}]
\textbf{Fase 3} --- Elasticidad, oscilaciones y calor

\textbf{Objetivo.} Verificar la ley de enfriamiento de Newton en dos cuerpos de distinto tamaño, determinar la constante de enfriamiento a partir de una linealización y comparar el tiempo característico de ambos.

\textbf{Contenidos}

\begin{itemize}[leftmargin=1.4em,itemsep=2pt,topsep=3pt]
  \item Ley de enfriamiento de Newton y régimen newtoniano.
  \item Diferencia de temperatura respecto del ambiente.
  \item Linealización logarítmica y constante de enfriamiento.
  \item Constante de tiempo y su relación con la geometría del cuerpo.
  \item Pérdidas de calor, convección y radiación.
  \item Registro temporal con termómetro digital y cronómetro.
\end{itemize}

\textbf{Herramientas:} Dos sólidos del mismo material y distinto tamaño, Termómetro digital con resolución de 0,1 $^{\circ}$C, Cronómetro digital, Plancha o baño de agua caliente, Pinzas y guantes térmicos, Base aislante y recipiente con agua.

\textbf{Práctica.} \emph{Enfriamiento de un cuerpo sólido}. Cada grupo calienta dos cuerpos del mismo material y distinto tamaño en agua, los retira y registra su temperatura cada 30 s durante 10 min. Debe linealizar la ley de enfriamiento, obtener la constante de cada cuerpo y explicar por qué el más pequeño se enfría más rápido.

\textbf{Entregable:} Informe final con las tablas de ambos cuerpos, las gráficas linealizadas, las constantes de enfriamiento con su incertidumbre, la comparación de los tiempos característicos y la discusión de las limitaciones del modelo.

\textbf{Instrumento de evaluación:} Práctica con informe y cierre del curso.
\end{tcolorbox}
\end{document}
